% -*- program: xelatex -*-

\newif\ifanswers
%\answerstrue

\documentclass[11pt]{article}

\usepackage{amssymb,amsmath,amsfonts,amsthm,graphicx}
\usepackage[top=0.3in, bottom=0.75in, left=0.5in, right=0.5in]{geometry}
\usepackage{tikz}
\usepackage{array}
\usepackage{enumitem}

\DeclareGraphicsExtensions{.pdf}

\usepackage{fontspec,realscripts}
\defaultfontfeatures{Ligatures=TeX}
\frenchspacing

\usepackage{xcolor}
\usepackage{pagecolor}
\usepackage[document]{ragged2e}

\begin{document}

\noindent
\begin{tabular*}{\textwidth}{@{\extracolsep{\fill}} l r}
\textbf{CS 310: Algorithms -- Matrix Chain Multiplication Worksheet Quiz}
& \emph{Mudassir Shabbir, March 25, 2026}
\end{tabular*}

\vspace{0.4cm}

\noindent
\begin{tabular*}{\textwidth}{@{\extracolsep{\fill}} l l}
\textbf{Student ID 1 \& 2:}{\em (Do this in pairs)} \hrulefill
&
\end{tabular*}

\vspace{0.3cm}
\hrule

\vspace{0.4cm}

\noindent
Matrix Chain rules for this worksheet:
\begin{itemize}
    \item Matrix dimensions are given by $p[0..n]$ for $n$ matrices.
    \item Matrix $A_i$ has size $p[i-1] \times p[i]$.
    \item Goal: choose a parenthesization that minimizes scalar multiplications.
\end{itemize}

\pagenumbering{gobble}

\begin{enumerate}[leftmargin=0pt, itemindent=1.7em, itemsep=1.2em]

\item \textbf{Warm-up: compare two parenthesizations}

For each dimensions array, compute both candidate costs and circle the cheaper order.

\vspace{0.2cm}
\begin{center}
\renewcommand{\arraystretch}{1.5}
\begin{tabular}{|>{\raggedright\arraybackslash}m{1.8in}|>{\raggedright\arraybackslash}m{2.2in}|>{\raggedright\arraybackslash}m{2.0in}|}
\hline
Dimensions $p$ & Candidate A & Candidate B \\\hline
$[10,30,5,60]$ & $(A_1A_2)A_3$ & $A_1(A_2A_3)$ \\\hline
$[5,10,3,12]$ & $(A_1A_2)A_3$ & $A_1(A_2A_3)$ \\\hline
$[8,20,2,25]$ & $(A_1A_2)A_3$ & $A_1(A_2A_3)$ \\\hline
\end{tabular}
\end{center}

\vspace{0.8cm}

\ifanswers
\textbf{Sample answers:}
\begin{itemize}
    \item $[10,30,5,60]$: $4500$ vs $27000$, so $(A_1A_2)A_3$.
    \item $[5,10,3,12]$: $330$ vs $960$, so $(A_1A_2)A_3$.
    \item $[8,20,2,25]$: $1320$ vs $4400$, so $(A_1A_2)A_3$.
\end{itemize}
\fi

\item \textbf{Fill the DP table}

Use $p=[10,30,5,60]$ and fill upper-triangular values for $\mathrm{MC}[i][j]$.

\vspace{0.2cm}
\begin{center}
\renewcommand{\arraystretch}{1.3}
\begin{tabular}{c|ccc}
  & $j=1$ & $j=2$ & $j=3$ \\\hline
$i=1$ & & & \\
$i=2$ & & & \\
$i=3$ & & & \\
\end{tabular}
\end{center}

\vspace{0.2cm}
What final value do you get for $\mathrm{MC}[1][3]$?

\vspace{1.0cm}

\ifanswers
\textbf{Sample filled table:}
\[
\begin{array}{c|ccc}
 & j=1 & j=2 & j=3 \\\hline
i=1 & 0 & 1500 & 4500 \\
i=2 &  & 0 & 9000 \\
i=3 &  &  & 0
\end{array}
\]
$\mathrm{MC}[1][3]=4500$.
\fi

\item \textbf{Choose the DP state}

Let there be $n$ matrices.

\begin{enumerate}[label=(\alph*), leftmargin=2em, itemsep=0.9em]
\item Complete the state definition:

\vspace{0.2cm}
\hspace*{1em}$\mathrm{MC}(i,j) =$ \hrulefill

\vspace{0.7cm}
\item What is the base case $\mathrm{MC}(i,i)$ and why?

\vspace{1.0cm}
\item Why do we fill by increasing chain length $\ell$?

\vspace{1.1cm}
\end{enumerate}

\ifanswers
\textbf{Sample answers:}
\begin{itemize}
    \item $\mathrm{MC}(i,j)$ is the minimum number of scalar multiplications to compute $A_iA_{i+1}\cdots A_j$.
    \item $\mathrm{MC}(i,i)=0$ because a single matrix needs no multiplication.
    \item Each interval $(i,j)$ depends on shorter intervals $(i,k)$ and $(k+1,j)$.
\end{itemize}
\fi

\newpage

\item \textbf{Build the recurrence and analyze complexity}

\begin{enumerate}[label=(\alph*), leftmargin=2em, itemsep=0.9em]
\item Write the recurrence using split point $k$.

\vspace{1.3cm}

\item Explain in one sentence what the extra term $p[i-1]p[k]p[j]$ represents.

\vspace{1.2cm}

\item What are the overall time and space complexities?

\vspace{1.0cm}
\end{enumerate}

\ifanswers
\textbf{Sample answers:}
\begin{itemize}
    \item
    \[
    \mathrm{MC}(i,j)=\min_{i\le k<j}\{\mathrm{MC}(i,k)+\mathrm{MC}(k+1,j)+p[i-1]p[k]p[j]\}.
    \]
    \item It is the cost of multiplying the two resulting matrices after solving the left and right subchains.
    \item Time: $O(n^3)$, Space: $O(n^2)$.
\end{itemize}
\fi

\end{enumerate}

\end{document}
